% Physical unit 015; accepted Session 02.
\GordanUnitStart{015}{printed page 10}
\GordanTopFolio{10}
In this series the constant coefficients of the individual powers of the function \gq{$\cos l_0$} are brought to simpler forms by putting:
\[
 p=o-{3\over2}o^2+{15\over4}o^3-{175\over8}o^4,
\]
or, by the inverse series,
\[
 o=p+{3\over2}p^2+{3\over4}p^3-{5\over16}p^4.
\]
When this series is substituted for $o$ in the equation, there results:
\[
 {2\mathcal A(l_0)\over\pi}=1-{b\over a}\{2p\cos l_0+(1-2p+{25\over8}p^2)p^2\cos^3l_0+({3\over2}-{51\over8}p)p^3\cos^5l_0+{25\over8}p^4\cos^7l_0\ldots\}.
\]

\begin{center}\S\ 5.\end{center}

Just as the definition of the geodetic line is sought from the period of extension, so it is sought from the period of longitude.

Accordingly, from $\mathcal A(l_0)$ the value $l_0$ can be computed explicitly; putting
\[
 a\left\{1-{2\mathcal A(l_0)\over\pi}\right\}=2bv\sqrt p\qquad \text{and}\qquad \cos l_0\sqrt p=x,
\]
one has:
\[
 v=x+x^3\left({1\over2}-p+{25\over16}p^2\right)+x^5\left({3\over4}-{51\over16}p\right)+{25\over16}x^7\ldots
\]
and, if we invert this equation,
\[
 x=v+v^3\left(-{1\over2}+p-{25\over16}p^2\right)+v^5\left({3\over16}p-{33\over32}p^2\right)-{1\over16}v^7.
\]

\begin{center}\S\ 6.\end{center}

Using the numbers explained in the first chapter, the forms are these:
\[
 o=0{,}001\ 6798\ 046\ 65\qquad p=0{,}001\ 1755\ 897\ 26
\]
\[
 \Sigma(l_0)={b\pi\over2}+{b\pi\over2\cdot10^{12}}\{167\ 9804\ 665\sin^2l_0-211\ 6291\sin^4l_0+5920\sin^6l_0-20\sin^8l_0\ldots\}
\]
\[
 \mathcal A(l_0)={\pi\over2}-{b\pi\over2\cdot10^{12}}\{335\ 1179\ 452\cos l_0+282\ 6309\cos^3l_0+7125\cos^5l_0+22\cos^7l_0\ldots\}.
\]

\begin{center}
{\large CHAPTER IV.}\\[0.8em]
\S\ 1.
\end{center}

The integral
\[
 \int_{l_1}^{l_2}dl\sqrt{1-e^2\cos^2l+\cos^2l\,{d\psi\over dl}\cdot{d\psi\over dl}}
\]
which, as we taught above, signifies the extension between the points $A(l_1,\psi_1)$ and $B(l_2,\psi_2)$ of any line whatsoever on the spheroid, becomes a minimum in a certain line.
